The update of Proposition \ref{prop.lobUpdate} is exact, and it is too fine for what we want.
It tracks the whole aggregated book, and to run it we need the whole tape of orders.
We look for something simpler:
a statistic computed from the observed aggregated book alone,
whose only job is to track the mid-price.

The order flow imbalance is introduced in \cite{CKS14pri} through a book in which the
mechanism of Section \ref{sec.orderDrivenMarkets} is simplified until the mid-price becomes
an explicit function of the flow.
We state the simplification as an assumption, and the consequence as a proposition.

\begin{assumption}\label{assumption.idealisedBook}
 Let $\uniformDepth > 0$ be fixed over the trading epoch, and assume that
 \begin{enumerate}[label={\emph{\roman{*}}.}]
  \item no event occurs behind the best quote of its side;
  \item no order reaches past the quote it acts on;
  \item every queue holds at most $\uniformDepth$,
        and every queue behind the best bid and the best ask holds exactly
        $\uniformDepth$;
  \item once the size at the best bid (ask, resp.) reaches $\uniformDepth$,
        the next bid (ask, resp.) limit order arrives one tick above (below, resp.)
        the best quote.
 \end{enumerate}
\end{assumption}

Parts \emph{i.}, \emph{iii.} and \emph{iv.} are the stylised book of \cite{CKS14pri}.
Part \emph{ii.} is not.
There the flows are totalled at whichever quote is best when each order arrives,
so an order reaching past the touch is split across the levels it crosses;
we ask instead that the contribution of an event be the size of that event,
which is Proposition \ref{prop.signedSizeContribution},
and part \emph{ii.} is what that costs.
The cap in \emph{iii.} is what \emph{iv.} intends,
and \eqref{eq.idealisedResidual} names a residual that the reference carries as the
truncation of a floor.

\begin{defi}[{\citealp{CKS14pri}}]\label{def.orderFlowImbalance}
 Let $\arrivalTimes[]_n$ be the times of the events of a trading epoch,
 and write $\nthBestBidSize[1]_{\arrivalTimes[]_n}$ for the size at the best bid
 \emph{after} the $n$-th event, so that
 $\nthBestBidSize[1]_{\arrivalTimes[]_{n-1}}
 = \nthBestBidSize[1]_{\arrivalTimes[]_n -}$.
 The contribution of the $n$-th event, for $n \geq 1$, is
 \begin{equation}
 \label{eq.orderFlowContribution}
 \begin{split}
  \orderFlowContribution
  =\;&
  \indicator{\bestBidPrice_{\arrivalTimes[]_n} \geq \bestBidPrice_{\arrivalTimes[]_{n-1}}}
  \nthBestBidSize[1]_{\arrivalTimes[]_n}
  -
  \indicator{\bestBidPrice_{\arrivalTimes[]_n} \leq \bestBidPrice_{\arrivalTimes[]_{n-1}}}
  \nthBestBidSize[1]_{\arrivalTimes[]_{n-1}}
  \\
  &-
  \indicator{\bestAskPrice_{\arrivalTimes[]_n} \leq \bestAskPrice_{\arrivalTimes[]_{n-1}}}
  \nthBestAskSize[1]_{\arrivalTimes[]_n}
  +
  \indicator{\bestAskPrice_{\arrivalTimes[]_n} \geq \bestAskPrice_{\arrivalTimes[]_{n-1}}}
  \nthBestAskSize[1]_{\arrivalTimes[]_{n-1}} ,
 \end{split}
 \end{equation}
 and the \emph{order flow imbalance} over the window $(t-\ofiWindow, t]$ is
 \begin{equation}
 \label{eq.orderFlowImbalance}
  \OFI_{t,\ofiWindow}
  = \sum_{n \, : \, t-\ofiWindow < \arrivalTimes[]_n \leq t} \orderFlowContribution .
 \end{equation}
\end{defi}

Both bid indicators fire when the best bid does not move,
so an event that leaves it in place contributes the change in the size resting there.
An event that moves it contributes a whole queue.

\begin{prop}\label{prop.idealisedUpdate}
 Under Assumption \ref{assumption.idealisedBook},
 the mid-price change over the window $(t-\ofiWindow,t]$ satisfies
 \begin{equation}\label{eq.idealisedUpdate}
  \Delta \midPrice_{t,\ofiWindow}
  = \midPrice_{t} - \midPrice_{t-\ofiWindow}
  = \frac{\tickSizeOfLOB}{2\uniformDepth} \, \OFI_{t,\ofiWindow}
  \; + \; \residual_{t,\ofiWindow},
  \qquad
  \vert \residual_{t,\ofiWindow} \vert \leq \tickSizeOfLOB .
 \end{equation}
 The residual has the explicit expression
 \begin{equation}\label{eq.idealisedResidual}
  \residual_{t,\ofiWindow}
  = -\frac{\tickSizeOfLOB}{2\uniformDepth}\Big[
  \big( \nthBestBidSize[1]_{t} - \nthBestBidSize[1]_{t-\ofiWindow} \big)
  - \big( \nthBestAskSize[1]_{t} - \nthBestAskSize[1]_{t-\ofiWindow} \big) \Big],
 \end{equation}
 and depends only on the sizes resting at the two touches at the endpoints of the window.
 In particular, \eqref{eq.idealisedUpdate} is exact whenever the bid and the ask touches
 change by the same amount over the window, and so whenever all four sizes equal
 $\uniformDepth$.
\end{prop}
\begin{proof}
 Proposition \ref{prop.lobUpdate} specialised to
 Assumption \ref{assumption.idealisedBook}; Exercise \ref{exercise.idealisedUpdate}.
\end{proof}

\begin{exercise}\label{exercise.idealisedUpdate}
 Prove Proposition \ref{prop.idealisedUpdate}.
 Measure the position of the bid in shares by
 $\nthBestBidSize[1]_{s}
 + \uniformDepth\,(\bestBidPrice_{s} - \bestBidPrice_{t-\ofiWindow})/\tickSizeOfLOB$,
 show that each event moves it by its own signed bid-side flow,
 and sum over the window.
 Say where each part of Assumption \ref{assumption.idealisedBook} is used.
\end{exercise}

\begin{remark}\label{remark.residualIsQueueFlow}
 On a window over which neither quote moves,
 $\Delta\midPrice_{t,\ofiWindow} = 0$ while
 $\OFI_{t,\ofiWindow}
 = \big(\nthBestBidSize[1]_{t} - \nthBestBidSize[1]_{t-\ofiWindow}\big)
 - \big(\nthBestAskSize[1]_{t} - \nthBestAskSize[1]_{t-\ofiWindow}\big)$,
 so \eqref{eq.idealisedResidual} gives
 $\residual_{t,\ofiWindow}
 = -\tickSizeOfLOB\,\OFI_{t,\ofiWindow}/(2\uniformDepth)$
 and the two terms of \eqref{eq.idealisedUpdate} cancel.
 The residual is the part of the flow that went into the queues rather than into the quotes.
 It is bounded by one tick and does not grow with $\ofiWindow$,
 which is what makes \eqref{eq.idealisedUpdate} useful over windows
 whose flow term is larger than that.
\end{remark}
